\documentclass[11pt,a4paper]{article}
\usepackage[margin=25mm]{geometry}
\usepackage{amsmath,amssymb,hyperref}
\newcommand{\contributorpreview}{\underline{\hspace{45mm}}}
\title{A distributed, verifiable search for $x^3+y^3+z^3=114$}
\author{Kuber Mehta\\math-gambling collaboration\\\small Local contribution preview: \contributorpreview}
\date{Working draft, September 2026 --- no result claimed}
\begin{document}
\maketitle
\begin{abstract}
We describe an open volunteer search for integer solutions to $x^3+y^3+z^3=114$. The implementation combines a cubic-field norm construction, exact arithmetic, modular rejection, bounded client tasks, durable receipts and independent replay. Scheduling adapts to measured execution cost while preserving exploration. The selected domains are finite, not an exhaustive height search. No solution is claimed in this draft. Contributor previews are not authorship or discovery claims.
\end{abstract}
\section{Problem and motivation}
The congruence obstruction modulo $9$ excludes targets congruent to $\pm4$. The target $114$ passes this necessary test, which does not prove that it has an integer solution. This project opens a previously local computational experiment to volunteer contributions.
\section{Search construction}
For $S=x+y$ and $V=x-y$,
\begin{equation}3SV^2=4(114-z^3)-S^3.\end{equation}
Writing $D=|S|$, a necessary condition is $r^3\equiv114\pmod D$, with $z=r+Dq$. Let $\alpha^3=114$. We use selected generators with
\begin{equation}N(a+b\alpha+c\alpha^2)=a^3+114b^3+12996c^3-342abc.\end{equation}
The task protocol fixes class lattices, coefficient regions and quotient bands. Congruence filters precede exact square, parity and integrality checks. Every emitted triple is checked against the original integer equation.
\section{Volunteer execution and independent verification}
Browser clients use JavaScript BigInt in an opt-in worker. The local runner uses Python arbitrary-precision integers with bounded process parallelism. Finished tasks have canonical receipts, which participants bank through GitHub issues. The issue author supplies authenticated submission identity. A scheduled Python verifier replays bounded tasks and deduplicates canonical identifiers. Submitted counters and claimed CPU time are not trusted. A SHA-256 digest binds a result but does not prove that the submitter ran a processor.
\section{Adaptive allocation}
After every 64 newly verified unique tasks, the cluster publishes a frozen cost policy. At least 40\% of allocation remains exploration. Measured replay efficiency informs the remainder. This is a throughput proxy, not an estimate of discovery probability. A separate exploratory discovery-learning experiment failed its acceptance criterion.
\section{Validation and limitations}
The archived native release recovered all 662 known-positive regression fixtures and included differential, arithmetic-boundary and fault-recovery tests. Those results are not automatically evidence for the portable worker, which has a separate cross-language suite. No formal proof of the complete software stack is claimed. Full replay of negative tasks duplicates computation and can bottleneck the verifier; linear cluster speedup has not been demonstrated. Unverified banks remain pending. Counter totals are not independent chances of finding a solution.
\section{Result --- intentionally unfilled}
\begin{description}
\item[Verified integer triple:] \rule{85mm}{0.2pt}
\item[Finder / verified contributor:] \rule{70mm}{0.2pt}
\item[Discovery date and receipt:] \rule{70mm}{0.2pt}
\item[Independent reproduction:] \rule{70mm}{0.2pt}
\end{description}
This section remains blank until independently verified evidence exists. A local name preview does not fill it.
\section{Reproducibility and attribution}
Source, fixed task definitions, tests, calibration versions and research notes are maintained at \url{https://github.com/Kuberwastaken/math-gambling}. Final discovery attribution and authorship require evidence and agreement; participation alone does not confer authorship.
\begin{thebibliography}{9}
\bibitem{booker} A. R. Booker and A. V. Sutherland. \emph{On a question of Mordell}. PNAS 118 (2021). \url{https://arxiv.org/abs/2007.01209}.
\bibitem{grantham} J. Grantham and P. G. Walsh. \emph{Representing integers as a sum of three cubes}. (2022). \url{https://arxiv.org/abs/2211.12149}.
\bibitem{huisman} S. G. Huisman. \emph{Newer sums of three cubes}. (2016). \url{https://arxiv.org/abs/1604.07746}.
\end{thebibliography}
\end{document}
